\section{Robot Localisation}

\subsection{Introduction}

Robot localisation is the process of estimating a robot's pose --- position and orientation --- within a given reference frame using sensory information. It is a prerequisite for all navigation and manipulation tasks that require interaction with the physical environment, and is tightly coupled with mapping; together they constitute Simultaneous Localisation and Mapping (SLAM) \cite{ullahMobileRobotLocalization2024}\cite{leonardSimultaneousMapBuilding1991}.

This section traces the historical foundations of robot localisation, structured around three persistent challenges that have shaped the field and that remain directly relevant to the methods presented in this thesis. The first is the accumulation of pose error under dead reckoning, which motivates a probabilistic treatment of uncertainty. The second is the need to represent and fuse uncertain, heterogeneous sensor measurements into a coherent spatial model. The third is the tension between short-range alignment and global map consistency. The foundational contributions reviewed here --- probabilistic uncertainty modelling \cite{smithRepresentationEstimationSpatial1987}, occupancy-based sensor fusion \cite{elfesUsingOccupancyGrids1989}, the Iterative Closest Point (ICP) algorithm \cite{beslMethodRegistration3D1992a}, and pose graph optimisation \cite{luGloballyConsistentRange1997} --- constitute the algorithmic backbone of the SLAM system developed in \cref{chap:SLAM} and the precise-alignment pipeline described in \cref{chap:precise-alignment}.

\subsection{The Limits of Deterministic Localisation}

The first mobile robots capable of autonomous reasoning operated under a deterministic model of the world: each commanded action was assumed to execute perfectly, and the robot's estimated pose was updated accordingly. Shakey the Robot, developed at Stanford Research Institute between 1966 and 1972, exemplified this approach \cite{nilssonShakeyRobot1984}. Shakey maintained no continuous pose estimate; localisation was expressed through discrete semantic primitives --- ``room 1'', ``doorway 1'' --- and updated by open-loop dead reckoning after each atomic motion command. Landmark-based correction using pre-modelled doorways and corners was available, but depended on the environment conforming exactly to the assumed structure, which required heavily controlled conditions to function at all. Nilsson's own assessment of the principal result was ``an appreciation of the potential complexity of the problem of vision when the real world is the subject matter'' \cite{nilssonShakeyRobot1984} --- a constraint that remains pertinent in unstructured outdoor environments.

Hans Moravec's Stanford AI Lab Cart extended the approach to stereo-vision-based geometric localisation \cite{moravecObstacleAvoidanceNavigatio1980}. The Cart extracted image features, estimated their 3D positions through stereo disparity, and derived its own motion by computing the rigid transformation between successive feature observations. Unlike Shakey's symbolic representation, the Cart maintained an explicit geometric model. However, motion estimation remained deterministic: no uncertainty was associated with feature positions or the recovered transformation, so localisation errors accumulated silently with no mechanism to detect or correct drift. The consequence was that for a 20-metre indoor course, navigation required approximately five hours \cite{moravecStanfordCartCMU1990}.

These two systems establish the core failure mode of deterministic localisation: uncorrected motion errors compound over time, and landmark-based correction is unreliable whenever the observed environment departs from assumed models. In semi-structured outdoor environments such as vineyards --- where plant geometry changes seasonally, features repeat along rows, and no guaranteed prior map exists --- this fragility is not a boundary case but the normal operating condition. The response was to replace deterministic pose estimates with probabilistic representations that could explicitly track and propagate uncertainty.

\subsection{Uncertainty Representation and Sensor Fusion}

The shift to probabilistic localisation was formalised by Smith and Cheeseman in ``On the Representation and Estimation of Spatial Uncertainty'' \cite{smithRepresentationEstimationSpatial1987}. They proposed representing each spatial relationship as an \textit{approximate transformation} (AT): a pose estimate paired with a covariance matrix encoding second-order uncertainty. Two operations on ATs underpin all subsequent probabilistic SLAM work.

\textit{Compounding} propagates uncertainty through sequential spatial transformations. Two uncertain poses are composed as
\[\mathbf{x}_3 = f(\mathbf{x}_1, \mathbf{x}_2),\]
where \(\mathbf{x}_i\) denotes a pose with mean \(\boldsymbol{\mu}_i\) and covariance \(\mathbf{C}_i\). Linearising \(f(\cdot)\) about the mean values via a first-order Taylor expansion gives
\[
\delta \mathbf{x}_3 \approx \mathbf{J}
\begin{bmatrix}
\delta \mathbf{x}_1 \\
\delta \mathbf{x}_2
\end{bmatrix},
\]
where \(\mathbf{J}\) is the Jacobian of \(f\) evaluated at \((\boldsymbol{\mu}_1, \boldsymbol{\mu}_2)\). The propagated covariance is then
\[
\mathbf{C}_3 = \mathbf{J}
\begin{bmatrix}
\mathbf{C}_1 & \mathbf{0} \\
\mathbf{0} & \mathbf{C}_2
\end{bmatrix}
\mathbf{J}^\top .
\]
This formalises what Shakey and the Cart demonstrated empirically: uncertainty grows with each compounded transformation, and correlations between position and orientation accumulate as a natural consequence of the transformation geometry.

\textit{Merging} reduces uncertainty by fusing independent estimates of the same pose. Given two estimates \(\mathbf{x}_1 \sim \mathcal{N}(\boldsymbol{\mu}_1, \mathbf{C}_1)\) and \(\mathbf{x}_2 \sim \mathcal{N}(\boldsymbol{\mu}_2, \mathbf{C}_2)\), the minimum-variance fused estimate is obtained via the Kalman gain
\[
\mathbf{K} = \mathbf{C}_1 \left( \mathbf{C}_1 + \mathbf{C}_2 \right)^{-1},
\]
with the updated mean and covariance
\[
\boldsymbol{\mu}_3 = \boldsymbol{\mu}_1 + \mathbf{K}\left( \boldsymbol{\mu}_2 - \boldsymbol{\mu}_1 \right),
\qquad
\mathbf{C}_3 = \mathbf{C}_1 - \mathbf{K}\mathbf{C}_1 .
\]
The merged covariance \(\mathbf{C}_3\) is smaller than either input, quantifying the information gain from combining independent sources. In a subsequent paper, Smith, Self, and Cheeseman applied this framework directly to SLAM through the Stochastic Map, in which sensor and landmark poses are jointly maintained in a single state vector with a shared covariance matrix --- the formulation that underpins Extended Kalman Filter (EKF) approaches to SLAM, including earlier iterations of the platform that motivates \cref{chap:SLAM}.

Building on this probabilistic foundation, Elfes introduced occupancy grids as a spatial representation capable of accommodating sensor uncertainty directly \cite{elfesUsingOccupancyGrids1989}. Where geometric maps required reliable feature extraction, occupancy grids maintain a probabilistic estimate of occupancy at each cell in a tessellated lattice, updated incrementally from raw sensor measurements through an explicit probabilistic sensor model. This allowed maps to be constructed from noisy, heterogeneous data without the brittleness of feature-based methods --- a critical property for environments with little or no prior geometric regularity.

A key insight from Elfes's work is that individual sensor modalities carry complementary uncertainty profiles: sonar provides low range uncertainty but poor angular resolution, whereas stereo cameras offer good angular resolution but uncertainty that grows with distance and degrades on textureless surfaces. Representing both in a common occupancy grid format allows them to be fused into a map with lower aggregate uncertainty than either source alone, as illustrated in \cref{fig:occupancyGridSensorFusion,fig:elfesMappingFramework}. Elfes also noted that fusing sensor views into a coherent global map requires accurate motion estimation, and that dead reckoning alone is insufficient precisely because uncertainty compounds in the manner formalised by Smith and Cheeseman. This co-dependence between localisation accuracy and map quality is the defining characteristic of SLAM, and directly motivates the choice in \cref{chap:SLAM} to fuse LiDAR odometry with an inertial measurement unit (IMU), whose complementary error characteristics reduce drift under the same merging principle.

\begin{figure}
    \centering
    \includegraphics[width=0.5\linewidth]{2-background/figures/globalMappingOccupancyGrid.png}
    \caption{A framework for occupancy-grid-based robot mapping. \protect\cite{elfesUsingOccupancyGrids1989}}
    \label{fig:elfesMappingFramework}
\end{figure}

\begin{figure}
    \centering
    \includegraphics[width=0.5\linewidth]{2-background/figures/sonarOccupancyGrid.png}
    \caption{Sonar mapping along a corridor \protect\cite{elfesUsingOccupancyGrids1989}.}
    \label{fig:sonarMapping}
\end{figure}

\begin{figure}
    \centering
    \includegraphics[width=0.5\linewidth]{2-background/figures/OccupancyGridSensorFusion.png}
    \caption{Sensor integration of sonar and scanline stereo. Occupancy grids generated separately for sonar (a) and scanline stereo (b), as well as the integrated map (c) obtained through sensor fusion, are shown. Occupied regions are marked by shaded squares, empty areas by dots fading to white space, and unknown spaces by + marks \protect\cite{elfesUsingOccupancyGrids1989}. }
    \label{fig:occupancyGridSensorFusion}
\end{figure}

\subsection{Point Cloud Registration}

Relative pose estimation between sensor frames requires a method for aligning observations acquired at successive robot positions. Besl and McKay's Iterative Closest Point (ICP) algorithm \cite{beslMethodRegistration3D1992a} provided this capability for 3D point clouds. ICP frames registration as an optimisation problem: starting from an initial alignment, it alternates between assigning each source point to its nearest neighbour in the target cloud and solving in closed form for the rigid transformation --- rotation and translation --- that minimises the mean squared distance over all correspondence pairs. Convergence to a local minimum is guaranteed; however, the solution is sensitive to initial conditions and requires sufficient overlap between the two clouds to avoid converging to a false minimum.

Subsequent work extended the basic formulation in directions that proved critical for practical SLAM. Chen and Medioni introduced the point-to-plane error metric \cite{chenObjectModelingRegistration1991}, which substantially improves convergence on structured surfaces. Rusinkiewicz et al.\ presented a suite of computational improvements that reduced per-iteration cost sufficiently to support real-time operation \cite{rusinkiewiczEfficientVariantsICP2001}. Segal et al.'s Generalised-ICP \cite{segalGeneralizedICP2009} unified point-to-point, point-to-plane, and plane-to-plane correspondences into a single probabilistic framework, improving robustness in environments that lack strong planar structure --- such as the organic, repetitive geometry of vineyard rows.

ICP is the direct registration engine used in the precise-alignment pipeline described in \cref{chap:precise-alignment}, where a real-time point cloud of a physical vine is aligned against the SfM-derived reference geometry extracted from a 3D Gaussian Splatting model. The sensitivity of ICP to initialisation and partial overlap is a central practical constraint addressed in that chapter.

\subsection{Pose Graph Optimisation and Global Consistency}

Probabilistic uncertainty modelling and local scan matching together support accurate short-range pose estimation, but do not address a deeper limitation: errors accumulated along a long trajectory cannot be corrected retroactively by filtering methods. EKF-based SLAM integrates past measurements into a global state and then discards those measurements; over extended trajectories, linearisation errors accumulate and the filter becomes internally inconsistent, with different regions of the map updated independently and unable to be reconciled.

Lu and Milios addressed this by introducing pose graph optimisation \cite{luGloballyConsistentRange1997}. In a pose graph, each estimated robot pose is a node, and each relative spatial constraint --- derived from scan matching or odometry --- is an edge. The optimisation objective is to find the configuration of all nodes that best satisfies the full set of edge constraints simultaneously. Because the entire history of pose estimates is retained, a single loop-closure detection event can propagate corrections backward through the trajectory, recovering a globally consistent map. The fundamental contrast with filtering is that past states are not discarded: global consistency is achieved through a batch optimisation rather than a recursive update, at the cost of a more demanding solve.

For long traversals through vineyard rows, where cumulative drift along the row axis would otherwise render multi-row maps geometrically inconsistent, this distinction is operative. The SLAM system presented in \cref{chap:SLAM} employs pose graph optimisation as its back-end, enabling loop-closure corrections to enforce global consistency across traversals of multiple vineyard rows.

\subsection{Summary}

The history reviewed in this section identifies three challenges that persist in the vineyard localisation problem addressed by this thesis. First, pose error accumulates under dead reckoning and must be explicitly modelled and propagated --- demanding a probabilistic treatment of estimation. Second, no single sensor is sufficient for robust localisation; each modality carries characteristic errors, and heterogeneous sensor fusion governed by the merging operations formalised by Smith and Cheeseman is necessary to reduce aggregate uncertainty. Third, local consistency from scan matching and global consistency from loop-closure correction are distinct requirements, and satisfying both simultaneously requires a graph-based optimisation back-end. The systems developed in \cref{chap:SLAM} and \cref{chap:precise-alignment} address these three challenges in sequence. The following section examines how these challenges manifest within vineyard and orchard environments, where domain-specific constraints have motivated a distinct body of applied SLAM research.

